% =====================================================================
% Conclusion
% =====================================================================
\section{Conclusion}
\label{sec:conclusion}

We have presented a comparative analysis of four numerical
interpolation methods drawn from Chapters~18 and~24 of Chapra and
Canale---linear interpolation, natural cubic splines, $C^4$ quintic
Hermite splines, and cubic B-spline approximation---applied to the
problem of robot trajectory generation. Every method was implemented
directly from the textbook formulation, validated against reference
implementations and analytic test functions, and exercised on two
realistic robotics scenarios sweeping over waypoint density,
geometric character, and time spacing.

The data support eight findings (Section~\ref{sec:discussion}), of
which three are non-obvious and worth re-emphasising. \emph{First},
adding waypoints under chord-length parameterisation \emph{worsens}
smoothness metrics for spline methods, despite the popular intuition
that more data is always better. \emph{Second}, the quintic spline
under natural boundary conditions produces \emph{higher} peak
acceleration and RMS jerk than the cubic on both datasets---an
artefact entirely attributable to the boundary condition rather
than the method, since cubic-derived clamped boundary conditions
reduce peak acceleration by 33\% and RMS jerk by 64\%.
\emph{Third}, the cubic B-spline absorbs sharp input geometry much
more gracefully than the interpolating splines, at the cost of a
guaranteed 0.28--0.92\,m waypoint deviation in our scenarios.

\subsection{Recommendations}
\label{sec:recommendations}

Based on the analysis, we recommend the following decision rule for
practitioners choosing among these methods:

\begin{itemize}[leftmargin=*]
\item If waypoints are strict constraints and the application has
known end-velocities and end-accelerations (e.g.\ stitching together
pre-planned segments), use a \textbf{clamped quintic spline}; this
configuration combines exact waypoint adherence with $C^4$
continuity and the smoothest derivative norms in our experiments.
\item If waypoints are strict constraints but the system is expected
to start and stop at rest with no other prior information, use a
\textbf{natural cubic spline}; the natural-BC quintic should be
avoided because of the endpoint jerk spike documented in
Section~\ref{sec:results-bc-sens}.
\item If waypoints are soft guidelines and the path contains
intentionally sharp geometric features (e.g.\ obstacle-avoidance
clearances), a \textbf{cubic B-spline approximation} is preferred for
its substantially smoother derivative profile.
\item Linear interpolation is appropriate only as a baseline or for
debugging; it should not be deployed without an additional
post-processing smoothing step because of the velocity discontinuity
at every waypoint.
\end{itemize}

\subsection{Limitations}
\label{sec:limitations}

The study has several limitations that constrain the generality of
the conclusions.

The datasets are synthetic. Real industrial waypoint sequences come
from path planners with their own geometric biases, and the
distribution of inputs in deployed systems may differ from our
constructed mixed/sharp/gradual variants. Replicating the analysis
on a corpus of real recorded trajectories would strengthen the
external validity of the findings.

The metrics are kinematic only. We do not include actuator-saturation
limits, mechanical resonance frequencies, payload disturbance
models, or thermal limits, all of which are relevant to trajectory
quality in deployed systems and could be incorporated into a
weighted composite metric.

Only two boundary-condition variants of the quintic spline are
explored. A more thorough sensitivity study would sweep over a
range of clamped end-velocities and end-accelerations, or
investigate alternative boundary specifications such as periodic or
not-a-knot conditions.

\subsection{Future work}
\label{sec:future-work}

Several extensions of this work are natural follow-ups for
Deliverable~4 or beyond:

\begin{enumerate}[leftmargin=*]
\item \textbf{Optimised B-splines.} The B-spline approximation has
free parameters in the choice of knot vector and the placement of
control points; a constrained optimisation that minimises waypoint
deviation while preserving smoothness would close the
waypoint-fidelity gap.
\item \textbf{Time-optimal reparameterisation.} The current study
fixes the time parameterisation a priori. Methods such as
time-scaling under acceleration or jerk constraints
\cite{macfarlane2003jerk} would let each method run at its physical
limit, providing a fairer comparison across methods.
\item \textbf{High-dimensional configuration spaces.} The 2-DOF arm
is the simplest non-trivial manipulator. Repeating the analysis on
a 6-DOF or 7-DOF arm would test whether the findings generalise.
\item \textbf{Comparison with geometrically-aware methods.} Methods
such as PH curves and minimum-jerk trajectories
\cite{kyriakopoulos1988minimum} provide alternatives that the
classical interpolation methods compared here do not, and a head-to-
head comparison would situate our results within the broader
trajectory-planning literature.
\end{enumerate}

\subsection{Reproducibility}

All code, datasets, and results are available in the project
repository, including:

\begin{itemize}[leftmargin=*]
\item Python implementations of all four interpolation methods,
the centered finite-difference schemes from Chapter~28, and the
metric and visualisation utilities;
\item the canonical CSV waypoint datasets and the parametric variant
generators used to construct the experimental sweep;
\item an executable Jupyter notebook that regenerates every figure
and table in this paper from the raw data;
\item the complete experiment metric table
(\texttt{full\_experiment.csv}, 28 rows $\times$ 14 columns).
\end{itemize}

The full experiment runs in approximately 30\,ms on commodity
hardware, supporting independent re-derivation of every reported
result.
